\section{系统：任务、环境与简单通用方法}

上一章讲为什么语言是入口，本章进入 Agent 作为系统的核心。姚顺雨把自己的研究概括成两个问题：第一，怎样做有价值、和现实世界更相关的任务和环境；第二，怎样做简单但通用的方法。前者是任务线，后者是方法线。很多讨论只看方法线，忽略任务线，但 Agent 的能力往往被环境定义。

\begin{figure}[H]
\centering
\includegraphics[width=0.96\textwidth]{agent-two-core-questions.png}
\caption{Agent 研究的两个核心：有价值的任务/环境 + 简单通用的方法。自制概念图，依据 00:17:50--00:35:00 对谈内容整理。}
\end{figure}

\begin{knowledgebox}{读图：任务和方法必须共同进化}
如果任务太简单，方法看起来很强但没有现实意义；如果任务足够真实但方法太复杂，就很难泛化。Agent 研究要同时设计环境、奖励、工具和通用方法。
\end{knowledgebox}

\subsection{Agent 的三波兴衰：方法线和任务线}

上一节说任务和方法要共同进化，本节把这个判断放回 Agent 历史。Agent 是一个古老概念：能自我决策、与环境交互并优化奖励的系统，都可以被称为 Agent。访谈提醒我们，Agent 的历史不是一条直线，而是多次兴衰：古典 AI/RL Agent、强化学习环境、再到 LLM Agent。每一波都不只是方法变化，也依赖任务和环境是否足够合适。

\begin{figure}[H]
\centering
\includegraphics[width=0.96\textwidth]{agent-history-waves.png}
\caption{Agent 三波兴衰：方法线和任务线相辅相成，不能只看算法。自制概念图，依据 00:17:50--00:29:00 对谈内容整理。}
\end{figure}

\begin{knowledgebox}{术语消化：Agent 系统的基本构件}
\begin{tabularx}{\textwidth}{p{0.20\textwidth}p{0.32\textwidth}X}
\toprule
构件 & 含义 & 为什么重要 \\
\midrule
Environment & Agent 观察和行动的世界 & 决定任务是否真实、反馈是否有效。 \\
Policy & 根据观察选择动作的策略 & 可以由 LLM、代码、规划器或多模型系统实现。 \\
Reward & 对动作结果的评价信号 & 决定 Agent 学什么、探索什么。 \\
Memory & 长期状态和经验记录 & 支撑跨任务、跨会话和个性化。 \\
Tools & 外部 API、代码、浏览器、文件等能力 & 让 Agent 从语言进入可执行世界。 \\
\bottomrule
\end{tabularx}
\end{knowledgebox}

\begin{knowledgebox}{方法线 vs 任务线}
\begin{tabularx}{\textwidth}{p{0.22\textwidth}p{0.34\textwidth}X}
\toprule
线索 & 关注什么 & 容易忽略什么 \\
\midrule
方法线 & 算法、模型、规划、训练技巧 & 任务是否足够真实，环境是否能反馈。 \\
任务线 & benchmark、环境、工具、数据和评价 & 方法是否简单通用，能否迁移到新任务。 \\
系统线 & 方法和任务如何形成闭环 & 计算成本、失败恢复、长期记忆和安全。 \\
\bottomrule
\end{tabularx}
\end{knowledgebox}

这张表也解释了为什么许多 Agent demo 看起来惊艳，却难以成为稳定产品。Demo 往往只展示方法线：模型会规划、会调用工具、会写几步动作；真实产品还要处理任务线：环境是否可重复、反馈是否可判定、失败是否可恢复、用户是否愿意把真实目标交给它。姚顺雨强调任务和环境，正是在提醒 Agent 研究不能只追求“会做样子”。

\begin{importantbox}{Agent benchmark 的设计标准}
一个好的 Agent benchmark 至少要满足三点：任务足够真实，反馈足够明确，环境足够开放。太简单的任务测不出泛化；太封闭的环境训练不出真实行动；没有可靠反馈的任务又无法形成学习信号。
\end{importantbox}

\subsection{Code 是关键 affordance}

前面讨论 Agent 的构件，本节进入“行动接口”。访谈里有一个非常重要的判断：code 像人的手，是 AI 最重要的 affordance。Affordance 指环境给予行动者的行动可能性。对 AI 来说，代码、API、命令行、浏览器和文件系统让模型不只是说话，而是能够改变数字世界。

\begin{figure}[H]
\centering
\includegraphics[width=0.96\textwidth]{code-affordance.png}
\caption{Code 是关键 Affordance：代码让 AI 从语言进入数字世界行动。自制概念图，依据 00:29:49--00:33:00 对谈内容整理。}
\end{figure}

\begin{importantbox}{为什么 code 是手}
语言表达意图，code 执行动作。只会语言的模型像“会想会说的人”；会写代码、调用工具和修改环境的 Agent，才开始拥有数字世界里的手。
\end{importantbox}

\begin{knowledgebox}{术语消化：Affordance}
Affordance 原本来自生态心理学，指环境向行动者提供的行动可能性。门把手给人“拉开门”的 affordance；按钮给人“点击”的 affordance；对 AI 来说，代码、API、命令行和浏览器提供“改变数字世界”的 affordance。
\end{knowledgebox}

\begin{knowledgebox}{数字世界 affordance 的层级}
\begin{tabularx}{\textwidth}{p{0.20\textwidth}p{0.34\textwidth}X}
\toprule
层级 & 例子 & Agent 能做什么 \\
\midrule
文本层 & 文档、网页、邮件 & 读取、摘要、改写、检索。 \\
代码层 & Python、SQL、shell、API & 计算、查询、调用服务、修改系统状态。 \\
界面层 & 浏览器、IDE、操作系统 & 跨工具执行任务，处理非结构化环境。 \\
组织层 & 多人协作、权限、审批 & 与人类流程共同完成高风险任务。 \\
\bottomrule
\end{tabularx}
\end{knowledgebox}

\subsection{本章小结}

Agent 系统的关键是任务、环境、方法和 affordance 的共同设计。语言模型提供认知能力，代码和工具提供行动能力，环境与 reward 提供学习方向。

